\chapter{Abstract}\label{abstract}

\textbf{Topic:} Development of a Deep Learning Model for Automated Tree Recognition and Green Space Mapping in Urban Environments.

\textbf{Authors:} Anuar Totin, Rasul Aidarkhanov, Berik Sharipov. Educational program: 6B06101 --- Computer Science (IT-2304). Scientific supervisor: Syndar Satbayev.

\textbf{Relevance.} Urban tree inventories are a basic input to municipal planning, environmental monitoring and climate-adaptation processes, but the manual field surveys that have historically supplied them are slow, expensive and quickly obsolete for a city of Astana's size. Modern deep-learning models for object detection and instance segmentation now operate on freely-available satellite imagery at the per-tree level, opening a route to an automated alternative --- but the published literature does not yet contain a single evaluation of these models on Central-Asian imagery.

\textbf{Aim.} To design, implement and evaluate an end-to-end software system that, given a satellite image of an area of Astana as input, automatically produces an inventory of the trees present in that area with per-tree polygon mask, confidence score and geographic coordinates.

\textbf{Object of research:} the urban green space of Astana, namely the population of individual trees observable from above in the visible spectrum. \textbf{Subject of research:} deep-learning models and software components for the detection, segmentation, geographic conversion and export of these trees from satellite imagery.

\textbf{Scientific novelty.} This is the first published evaluation of state-of-the-art deep-learning tree-detection models on Astana satellite imagery and, by extension, on any major Central-Asian capital. The work delivers (i) the first measured magnitude of the geographic-generalisation gap for the region --- Box mAP@50 = \textbf{0.012} for the public NEON-pretrained DeepForest checkpoint on Astana, three orders of magnitude below the same model's reported off-the-shelf F-score on NAIP USA imagery --- (ii) a hybrid four-model architecture combining YOLOv8-seg instance segmentation, Mask R-CNN instance segmentation, fine-tuned DeepForest bounding-box detection and zero-shot Segment Anything 2 (SAM 2) mask refinement of DeepForest boxes, combined through both a Weighted-Box-Fusion ensemble and a novel cross-YOLO voting ensemble, and (iii) a custom-annotated Astana dataset of approximately 100 source images and ≈ 8 700 polygon-level annotations spanning three iterative batches (v1 / v2 / v3) that is reusable for future regional research.

\textbf{Methods.} Literature analysis of 31 peer-reviewed publications of 2019 -- 2025; deep-learning training and fine-tuning in PyTorch / Ultralytics / torchvision / DeepForest / SAM 2; dataset engineering in CVAT with a model-in-the-loop pre-labelling tool reducing annotation cost by approximately 70 \% per image; quantitative evaluation through standard COCO-style object-detection metrics (Box mAP@50, mAP@50:95, Mask mAP@50, Mask mAP@50:95, precision, recall) on a single 14-image cross-model validation set (set M14, 702 polygons) ensuring apples-to-apples comparison between all branches; multi-replicate variance estimation across four independent training runs of the best configuration; visual cross-checkpoint qualitative comparison on representative Astana scenes.

\textbf{Headline results.} A structured \textbf{23-experiment hyperparameter ablation} across six orthogonal axes (model size, starting weights, optimiser, augmentation level, inference resolution, training-time data ordering) and a four-replicate variance check selected the production YOLO configuration. On the M14 cross-model validation set, the best single configuration --- a YOLOv8x-seg checkpoint trained from public COCO weights with the Ultralytics default augmentation pipeline --- reaches Box mAP@50 = \textbf{0.315} and Mask mAP@50 = \textbf{0.289}, a +140 \% relative improvement on box detection over the YOLOv8x-seg v1 baseline (0.131) and +69 \% over the previous v2-finetune production (0.187). A counter-intuitive secondary finding emerges from the size ablation: the ``default Ultralytics augmentation'' pipeline (aggressive HSV colour jitter, random erasing, no geometric transformations) \textbf{outperforms the manually-tuned augmentation} that produced every preceding YOLO checkpoint on this same dataset, suggesting that the satellite-tree distribution is best served by colour-only augmentation since trees viewed from above are rotation-invariant and geometric transformations inject unnatural variance. Mask R-CNN v2+v3 (Box 0.166 / Mask 0.158) and DeepForest v3 + SAM 2 (Box 0.146 / Mask 0.134) follow as the second and third best branches respectively. The integrated pipeline processes a 1 km × 1 km capture at zoom 19 in approximately 18 seconds on a single laptop GPU.

\textbf{Practical result.} A deployable prototype --- FastAPI backend with a pluggable four-model adapter interface (eight YOLO checkpoint variants exposed through a hierarchical model picker plus DeepForest, DeepForest + SAM 2, Mask R-CNN, WBF ensemble and cross-YOLO 4-way vote ensemble), React 18 + Leaflet frontend with dark-mode-default styling, four-mode geographic conversion, in-browser ESRI / Google tile capture with Auto-Zoom Region Scan and streaming NDJSON progress, SQLite persistent storage of every snapshot / run / detection ever produced by the system, and three exporters (GeoJSON, CSV, standalone HTML) --- that delivers an Astana tree inventory in a single click and serves as a reusable template for any other Central-Asian city. The system meets all functional requirements set by \emph{Zelenstroy} and establishes the empirical baseline against which any future model for the region will be compared.

\newpage
